\documentclass[10pt,a4paper]{article}
\usepackage{webclases-lesson}
\lessonSetup{T1}{v0.2}{T1 - Movimiento rectilineo}

\begin{document}
\shorthandoff{<>}
\color{Ink}

\coverHeader{T1 --- Movimiento rectilineo}{Modelos de movimiento en una dimension y lectura de graficas.}

\begin{lessoncard}{Proposito}
La cinematica no empieza memorizando formulas: empieza eligiendo un eje, un origen, un reloj y un modelo. En esta leccion distinguimos MRU, MRUA y lectura grafica sin repetir la misma relacion en cada pagina.
\begin{keybox}
\textbf{Clave de lectura.} El signo de $x$, $v$ o $a$ depende del eje elegido; la rapidez no lleva signo, la velocidad si.
\end{keybox}
\end{lessoncard}

\begin{lessoncard}{Indice por bloques}
\footnotesize
\begin{tabularx}{\linewidth}{@{}>{\raggedright\arraybackslash}p{1.6cm}Y>{\raggedleft\arraybackslash\bfseries}p{1.0cm}@{}}
\textbf{Bloque} & \textbf{Alcance} & \textbf{Pag.}\\
\midrule
\tocBlock{Bloque 1}{Describir movimiento en una dimension}
\tocLine{sec:b1}{1}{Sistema de referencia, posicion, desplazamiento, velocidad y signos}
\tocBlock{Bloque 2}{Velocidad constante}
\tocLine{sec:b2}{2}{MRU, graficas y encuentros entre dos cuerpos}
\tocBlock{Bloque 3}{Aceleracion constante}
\tocLine{sec:b3}{3}{MRUA, frenadas, cambios de sentido, vertical y parada}
\tocBlock{Bloque 4}{Metodo grafico y sintesis}
\tocLine{sec:b4}{4}{Pendientes, areas, distancia frente a desplazamiento y checklist}
\end{tabularx}
\end{lessoncard}

{\scriptsize\textcolor{Muted}{Fuentes: PDF original T1, background editorial T1 y contraste conceptual con OpenStax University Physics 2. Diagramas redibujados.}}

\clearpage

\lessonHeader{Bloque 1 \textperiodcentered\ Describir movimiento}{Apartados internos}{Eje, signos y magnitudes}{Objetivo: convertir una historia en variables con sentido.}
\begin{lessoncard}{Referencia y magnitudes}
\subtopic{1}{Describir movimiento en una dimension}{sec:b1}
Elegir un sistema de referencia significa fijar origen, eje positivo y reloj. Entonces:
\[
\Delta x=x_f-x_i,\qquad d=\text{longitud recorrida},\qquad
v_m=\frac{\Delta x}{\Delta t},\qquad rapidez_m=\frac{d}{\Delta t}.
\]
El desplazamiento puede ser cero aunque la distancia no lo sea. La aceleracion mide cambio de velocidad:
\[
a_m=\frac{\Delta v}{\Delta t}.
\]
\end{lessoncard}

\begin{center}
\begin{tikzpicture}[x=.9cm,y=.75cm]
  \draw[axis] (-4.5,0) -- (5,0) node[right] {\scriptsize $x$ (m)};
  \foreach \x in {-4,-2,0,2,4} \draw (\x,.07)--(\x,-.07) node[below] {\scriptsize \x};
  \draw[vector] (-3,.45) -- (2.5,.45) node[midway,above] {\scriptsize $\Delta x=+5,5\units{m}$};
  \draw[vectorgreen] (-3,-.45) -- (-.5,-.45);
  \draw[vectorgreen] (-.5,-.45) -- (2.5,-.45) node[midway,below] {\scriptsize distancia suma tramos};
  \fill[BrandBlue] (-3,0) circle (.08) node[above] {\scriptsize inicio};
  \fill[AmberLine] (2.5,0) circle (.08) node[above] {\scriptsize final};
\end{tikzpicture}
\end{center}

\begin{warnbox}
``Va hacia la izquierda'' no significa necesariamente que acelere hacia la izquierda. Mira si la velocidad se hace mas negativa, menos negativa o cambia de signo.
\end{warnbox}

\clearpage

\lessonHeader{Bloque 2 \textperiodcentered\ Velocidad constante}{Apartados internos}{MRU y encuentros}{Objetivo: reconocer cuando la pendiente de $x(t)$ es constante.}
\begin{lessoncard}{Modelo MRU}
\subtopic{2}{Velocidad constante}{sec:b2}
Si la velocidad es constante:
\[
x(t)=x_0+vt,\qquad a=0.
\]
La grafica $x$-$t$ es una recta; su pendiente es $v$. La grafica $v$-$t$ es horizontal; el area bajo ella es el desplazamiento.
\begin{center}
\begin{tikzpicture}[x=.75cm,y=.55cm]
  \draw[axis] (0,0) -- (5.2,0) node[right] {\scriptsize $t$};
  \draw[axis] (0,0) -- (0,4.0) node[above] {\scriptsize $x$};
  \draw[BrandBlue,line width=1pt] (.4,.8) -- (4.7,3.2);
  \node at (3.6,2.0) {\scriptsize pendiente $v$};
  \begin{scope}[xshift=7cm]
    \draw[axis] (0,0) -- (5.2,0) node[right] {\scriptsize $t$};
    \draw[axis] (0,0) -- (0,4.0) node[above] {\scriptsize $v$};
    \draw[BrandBlue,line width=1pt] (.4,2.2) -- (4.7,2.2);
    \fill[GreenLine!20] (.8,0) rectangle (3.8,2.2);
    \node at (2.3,1.1) {\scriptsize area $=\Delta x$};
  \end{scope}
\end{tikzpicture}
\end{center}
\end{lessoncard}

\begin{examplebox}{encuentro en el mismo eje}
Desde $x_A(0)=-20\units{m}$, A camina a $+2,0\mps$. Desde $x_B(0)=40\units{m}$, B camina a $-1,0\mps$. Halla encuentro.
\solutionblock
\[
x_A=-20+2t,\qquad x_B=40-t,\qquad x_A=x_B.
\]
\[
-20+2t=40-t \Rightarrow t=20\units{s},\qquad x=20\units{m}.
\]
La condicion de encuentro es misma posicion en el mismo instante, no misma distancia recorrida.
\end{examplebox}

\clearpage

\lessonHeader{Bloque 3 \textperiodcentered\ Aceleracion constante}{Apartados internos}{MRUA, frenadas y vertical}{Objetivo: usar el modelo solo cuando $a$ sea constante.}
\begin{lessoncard}{Modelo MRUA}
\subtopic{3}{Aceleracion constante}{sec:b3}
\[
v=v_0+at,\qquad x=x_0+v_0t+\frac12at^2,\qquad v^2=v_0^2+2a\Delta x.
\]
La aceleracion tiene el signo del cambio de velocidad. Si $v$ y $a$ tienen signos opuestos, la rapidez disminuye; si tienen el mismo signo, aumenta.
\end{lessoncard}

\begin{examplebox}{frenar y cambiar de sentido}
Un carrito parte de $x_0=100\units{m}$ con $v_0=12\mps$ y $a=-3,0\mpss$.
\solutionblock Se detiene cuando $0=12-3t$, luego $t=4,0\units{s}$. La posicion entonces es
\[
x=100+12(4,0)-\frac12 3,0(4,0)^2=124\units{m}.
\]
Despues de $4,0\units{s}$, si la misma aceleracion continua, la velocidad se vuelve negativa y el movil regresa.
\end{examplebox}

\begin{examplebox}{movimiento vertical}
Lanzas una pelota hacia arriba con $v_0=18\mps$ desde $y_0=1,5\units{m}$ y tomas $+y$ hacia arriba, $g=10\mpss$.
\solutionblock
\[
v=18-10t,\qquad y=1,5+18t-5t^2.
\]
Altura maxima: $v=0\Rightarrow t=1,8\units{s}$, $y=17,7\units{m}$. Al caer por el punto de lanzamiento lleva $v=-18\mps$ si despreciamos rozamiento.
\end{examplebox}

\clearpage

\lessonHeader{Bloque 3 \textperiodcentered\ Aceleracion constante}{Aplicacion}{Reaccion y distancia de parada}{Objetivo: separar tramos con modelos distintos.}
\begin{examplebox}{reaccion y frenado}
Un coche circula a $72\units{km/h}=20,0\mps$. La reaccion dura $0,80\units{s}$ y luego frena con $a=-5,0\mpss$.
\solutionblock Durante la reaccion hay MRU:
\[
d_r=20,0(0,80)=16,0\units{m}.
\]
Durante la frenada:
\[
0=v_0^2+2a\Delta x \Rightarrow \Delta x=\frac{-v_0^2}{2a}=\frac{-400}{-10}=40,0\units{m}.
\]
Distancia total: $56,0\units{m}$. Si duplicas la velocidad, el tramo de reaccion se duplica, pero el de frenada se cuadruplica.
\end{examplebox}

\begin{lessoncard}{Cuando no usar una ecuacion}
No mezcles tramos con aceleraciones distintas en una sola formula. Primero divide el movimiento: reaccion, frenada, subida, bajada, espera o cambio de modelo. Cada tramo necesita condiciones iniciales propias.
\end{lessoncard}

\clearpage

\lessonHeader{Bloque 3 \textperiodcentered\ Aceleracion constante}{Comparacion de modelos}{Tablas y graficas rapidas}{Objetivo: decidir MRU o MRUA sin que el enunciado lo diga.}
\begin{lessoncard}{Senales de modelo}
\begin{tabularx}{\linewidth}{@{}p{3.0cm}Y Y@{}}
\toprule
Datos observados & Modelo probable & Comprobacion\\
\midrule
$x$ cambia lo mismo cada segundo & MRU & pendiente $x$-$t$ constante\\
$v$ cambia lo mismo cada segundo & MRUA & pendiente $v$-$t$ constante\\
$a=0$ & MRU o reposo & $v$ constante, no necesariamente cero\\
$v=0$ en un instante & no decide & mira $a$ para saber que ocurre despues\\
\bottomrule
\end{tabularx}
\end{lessoncard}

\begin{examplebox}{tabla corta}
Una tabla da $t=0,1,2,3\units{s}$ y $v=2,5,8,11\mps$.
\solutionblock La velocidad aumenta $3\mps$ cada segundo: $a=3\mpss$ constante. Es MRUA. Si $x_0=0$, entonces
\[
x(t)=2t+\frac12(3)t^2.
\]
En $3\units{s}$: $x=6+13,5=19,5\units{m}$.
\end{examplebox}

\begin{warnbox}
No conviertas un punto donde $v=0$ en reposo permanente. En una frenada con aceleracion constante, ese punto puede ser solo el instante de cambio de sentido.
\end{warnbox}

\clearpage

\lessonHeader{Bloque 4 \textperiodcentered\ Metodo grafico y sintesis}{Apartados internos}{Pendientes, areas y decision}{Objetivo: resolver incluso cuando no hay una formula ya preparada.}
\begin{lessoncard}{Leer graficas}
\subtopic{4}{Metodo grafico y sintesis}{sec:b4}
\begin{tabularx}{\linewidth}{@{}p{3.2cm}Y@{}}
\toprule
Grafica & Lectura\\
\midrule
$x(t)$ & pendiente $=v$; curvatura indica cambio de velocidad\\
$v(t)$ & pendiente $=a$; area con signo $=\Delta x$\\
$a(t)$ & area con signo $=\Delta v$\\
\bottomrule
\end{tabularx}
\end{lessoncard}

\begin{examplebox}{$v(t)=4-2t$ entre $0$ y $4\units{s}$}
\begin{center}
\begin{tikzpicture}[x=1.0cm,y=.55cm]
  \draw[axis] (0,0) -- (5,0) node[right] {\scriptsize $t$ (s)};
  \draw[axis] (0,-2.4) -- (0,4.6) node[above] {\scriptsize $v$ (m/s)};
  \draw[fill=GreenLine!20,draw=none] (0,0)--(0,4)--(2,0)--cycle;
  \draw[fill=AmberLine!20,draw=none] (2,0)--(4,-4)--(4,0)--cycle;
  \draw[BrandBlue,line width=1pt] (0,4)--(4,-4);
  \node at (1.1,1.3) {\scriptsize $+4\units{m}$};
  \node at (3.1,-1.3) {\scriptsize $-4\units{m}$};
\end{tikzpicture}
\end{center}
\solutionblock El desplazamiento total es $0\units{m}$, pero la distancia recorrida es $8\units{m}$. El cambio de sentido ocurre cuando $v=0$, en $t=2\units{s}$.
\end{examplebox}

\begin{exercisebox}{reconocer MRUA en una tabla}
La tabla da $t$ en s: $0,1,2,3$ y $x$ en m: $10,12,18,28$. Comprueba si puede venir de un MRUA desde reposo.
\solutionblock Los incrementos de posicion son $2,6,10$: crecen de cuatro en cuatro, compatible con aceleracion constante. Con $x=x_0+\frac12at^2$ y $x_0=10$, usando $t=1$: $2=\frac12a$, luego $a=4\mpss$. Verifica $t=2$: $10+8=18$ y $t=3$: $10+18=28$.
\end{exercisebox}

\begin{lessoncard}{Checklist de modelo}
1. Eje, origen y signo positivo. 2. Datos con unidades SI. 3. Decide MRU, MRUA o grafica. 4. Escribe la condicion fisica: encuentro, parada, altura maxima, distancia. 5. Comprueba signo, unidades y orden de magnitud.
\end{lessoncard}

\end{document}
